% Shared exercise chunk: l4-discrepancy-reverse
\begin{exercisechunk}
\item \textbf{Reverse discrepancy conversion.\suppmark}
\label[exercise]{ex:discrepancy-reverse}
\exerciseprofile{2}{2}{2}{\exproof\quad\excalculation\quad\exconstruction\quad\exinterpretation}{%
Determine when small task cost controls forward KL and derive the sharp
fixed-cost interval.}
Use the definitions in \cref{sec:discrepancy-comparison}. Let
\(P=\delta_{0^T}\), \(T\ge2\), with
\(c_{\mathrm{cnt}}(x)=\sum_{t=1}^T x_t\) and
\(c_{\mathrm{fin}}(x)=x_T\).
\begin{enumerate}[(a),beginpenalty=10000]
\item Use the fixed-KL intervals in
\cref{eq:binary-kl-level,eq:terminal-kl-level} to calculate the reverse
conversions:
\begin{align}
 \omega_{c_{\mathrm{cnt}}\to D}(s;P)
 &=\begin{cases}-\log(1-s),&0\le s<1,\\
                 +\infty,&s\ge1,\end{cases}
 \label{eq:binary-converse-modulus}\\
 \omega_{c_{\mathrm{fin}}\to D}(s;P)
 &=+\infty,\qquad s\ge0.
 \label{eq:terminal-converse-modulus}
\end{align}
Exhibit laws attaining each finite bound and laws with infinite KL
where appropriate. Explain why small expected count forces small
KL, whereas perfect final answers allow arbitrarily large KL.
\item Fix an expected count \(s\in[0,1)\). Prove the sharp bounds
\[
 -\log(1-s/T)\le\KL(P, Q)\le-\log(1-s).
\]
Give a law attaining each endpoint and relate them to part (a).
\end{enumerate}
\end{exercisechunk}
